\documentclass[12pt]{article}
\usepackage[utf8]{inputenc}
\usepackage[T1]{fontenc}
\usepackage{hyperref}
\title{Dass Language Resources for K-12 Curriculum Development}
\author{Compiled by Claude Sonnet 5.5 (claude-sonnet-5-5) for Andr\'as Kornai}
\date{}
\begin{document}
\maketitle
\section{Introduction}
Dass (ISO 639-3 \texttt{dot}; Glottocode dass1243) is a West Chadic (Afro-Asiatic) dialect cluster of Bauchi and Plateau States, Nigeria, also called Barawa. Searches found very little: no school materials, primers or standard orthography were located.
\section{Primary Educational Resources}
None found in search results.
\section{Linguistic and Reference Materials}
\begin{itemize}
\item Wikipedia, ``Dass language'': overview of the language and its varieties. \url{https://en.wikipedia.org/wiki/Dass_language}
\item Cosper, Ronald (1999), \emph{Barawa Lexicon} (LINCOM): a multi-language wordlist of South Bauchi West Chadic languages that includes Dass (as ``Dott''); bibliographic record on Glottolog. \url{https://glottolog.org/resource/reference/id/136865}
\end{itemize}
\section{Religious and Community Resources}
None found in search results.
\section{Digital and Audio Resources}
None found in search results.
\section{Curriculum Development Context}
There is no evidence of a standard orthography or of teaching in the language. Any curriculum would first require orthography development and community consultation; the Cosper wordlist could serve as a starting lexical source.
\end{document}